\scanpage{48}

and a morphism of \deleted{\unsure{functors}} \inserted{presheaves} over $S_0$
\[
(52)\qquad
\diagram{\begin{tikzcd}[row sep=small]
\widetilde{\underline Z}_{\rho,\sigma} \arrow[r] & \underline{\operatorname{Aut}}_{\mathrm{gr}}\bigl(\unsure{\underline C^{+}|\underline G^{+}}\bigr),\\
\underline u \arrow[r,mapsto] & u
\end{tikzcd}}
\]
giving rise to a morphism
\[
(53)\qquad
\diagram{\begin{tikzcd}[row sep=small]
\widetilde{\underline Z}_{\rho,\sigma} \arrow[r] & \underline{\operatorname{Aut}}_{\mathrm{gr}}\bigl(\unsure{\underline C^{+}|\underline G^{+}}\bigr)\times\mathbf G_{\mathrm m}\times\mu_{S_0}\times\mu_{S_0}\\
\underline u \arrow[r,mapsto] & u,\mu(\underline u),\varepsilon_2(\underline u),\varepsilon_3(\underline u)
\end{tikzcd}}
\]
--- calculations already done thirty-six times in various contexts should show that this morphism (which is \uline{monomorphic}\ \unsure{?}\ in characteristic $\ne 2,3$) is a morphism of \uline{groups}, for a suitable \uline{\unsure{group|groups}} structure on $\widetilde{\underline Z}_{\rho,\sigma}$, uniquely determined (since it is so in characteristic $\ne 2,3$, and we conclude by reason of schematic density). We will have \deleted{\unsure{this is not}}
\[
(54)\qquad
\diagram{\begin{tikzcd}
1\arrow[r]&\widetilde{\underline Z}_{\rho,\sigma}^{+0}\arrow[r]&\widetilde{\underline Z}_{\rho,\sigma}\arrow[r]&\mathbf G_{\mathrm m}\times\mu_{S_0}\times\mu_{S_0}\arrow[r]&1
\end{tikzcd}}
\]
\margin{a bit \unsure{cavalier}, since
\diagram{\begin{tikzcd}
\underline Z_{\rho,\sigma}\arrow[r]&\mathbf G_{\mathrm m}\times\mu_{S_0}\times\mu_{S_0}
\end{tikzcd}}
is not an epimorphism, because
\diagram{\begin{tikzcd}[row sep=small]
\underline M\times\underline M\arrow[r]&\underline M\\
(\alpha_0,\beta_0)\arrow[r,mapsto]&(1-\rho_{\underline M})\alpha_0\unsure{-|+}(1-\sigma_{\underline M})\beta_0
\end{tikzcd}}
is not epimorphic --- \unsure{where there} is a \unsure{swindle} in characteristic 3 --- cf. below\ldots}
making explicit
\[
(55)\qquad
\begin{gathered}
\widetilde{\underline Z}_{\rho,\sigma}^{+0}\underset{\inserted{\text{subgroup scheme}}}{\subset}\underline M\times\underline M\times\mathbf G_{\mathrm a}\\
\Vert\\
\left\{(\alpha_0,\beta_0,\nu')\,\middle|\,
\begin{gathered}
\deleted{\unsure{\alpha_0}}4\nu'=0\\
(1+\rho_{\underline M})\alpha_0=(1-\sigma_{\underline M})\beta_0+2\nu'\lambda_1
\end{gathered}\right\}
\end{gathered}
\]
We will thus have
\[
(56)\qquad
\diagram{\begin{tikzcd}
1\arrow[r]&\widetilde{\underline Z}_{\rho,\sigma}^{+00}\arrow[r]&\widetilde{\underline Z}_{\rho,\sigma}^{+0}\arrow[r]&\underline{{}_4\mathbf G_{\mathrm a}}\arrow[r]&1
\end{tikzcd}}
\]
\margin{47}
\margin{47\\\texttt{https://grothendieck.umontpellier.fr}\\Université de Montpellier, UMR CNRS 5149, Case courrier 051, Place Eugène Bataillon\\34095 Montpellier cedex 5 - France}

\scanpage{49}
\mspage{672}

where
\[
(57)\qquad
\begin{gathered}
\widetilde{\underline Z}_{\rho,\sigma}^{+00}\overset{\mathrm{dfn}}{=}
\left\{\underline u\in\Gamma\widetilde{\underline Z}_{\rho,\sigma}\,\middle|\,
\nu'(\underline u)=\nu_2(\underline u)=\nu_3(\unsure{u|\underline u})=0\right\}\\
\simeq\\
\left\{(\alpha_0,\beta_0)\in\Gamma(\underline M\times\underline M)\,\middle|\,
(1-\rho_{\underline M})\alpha_0=(1-\sigma_{\underline M})\beta_0\right\}.
\end{gathered}
\]
\margin{in characteristic 3 its relative dimension \unsure{jumps} from 2 to 3 !}
\uline{group} scheme $\underline Z_{\rho,\sigma}^{+00}$ is smooth \inserted{in characteristic $\ne 3$} of relative dimension 2 --- it is even canonically a Module scheme, we will have
\[
(58)\qquad\widetilde{\underline Z}_{\rho,\sigma}^{+00}\simeq\unsure{\mathbf V|\mathbb V}(\underline P)
\]
where $P$ is the free $\mathbb Z$-module of rank 2 defined \deleted{\illegible{1}} by the exact sequence
\[
(59)\qquad
\diagram{\begin{tikzcd}[row sep=small]
0\arrow[r]&\deleted{\unsure{M}}\inserted{P}\arrow[r]&M\times M\arrow[r]&M\arrow[r]&0\\
&&(\alpha_0,\beta_0)\arrow[r,mapsto]&(1-\rho_{\underline M})\alpha_0-(1-\sigma_{\underline M})\beta_0&
\end{tikzcd}}
\]
The homomorphism $M\times M\longrightarrow M$ is made explicit, setting
\[
(60)\qquad
\left\{\begin{aligned}
\alpha_0&=a_0\lambda_0+\unsure{a_1}\lambda_1
&\text{hence}\quad\alpha_0-\rho\alpha_0&=(a_0+\unsure{a_1})\lambda_0+(2\unsure{a_1|b_1}-a_0)\lambda_1\\
\beta_0&=b_0\lambda_0+b_1\lambda_1
&\text{hence}\quad\beta_0-\sigma\beta_0&=(b_0-b_1)\lambda_0+(b_1-b_0)\lambda_1
\end{aligned}\right.
\]
as
\[
(61)\qquad
\diagram{\begin{tikzcd}
(\alpha_0,\beta_0)\arrow[r,mapsto]&
\underbrace{(a_0+a_1+(b_1\unsure{-|+}b_0))}_{c_0}\lambda_0+
\underbrace{(2a_1-a_0-(b_1-b_0))}_{c_1}\lambda_1
\end{tikzcd}}
\]
\deleted{\unsure{We}} This is a surjective homomorphism \inserted{except in characteristic 3}, \deleted{for given $(c_0,c_1)$,}
\deleted{\unsure{we will set}\ $\unsure{a_1|c_1}$\ \illegible{1}}
\inserted{--- the elements of the image are the}
\[
(62)\qquad
\left\{c_0\lambda_0+c_1\lambda_1\,\middle|\,
\underbrace{c_0+c_1}_{=3a_1}\equiv0\quad(3)\right\}
\]
\uline{Remark} If we take \deleted{\unsure{a}}\inserted{a} $u$ defined over $\widehat{\mathbb Z}$ coming from an element of $M[0]$, then
\margin{48}
\margin{48\\\texttt{https://grothendieck.umontpellier.fr}\\Université de Montpellier, UMR CNRS 5149, Case courrier 051, Place Eugène Bataillon\\34095 Montpellier cedex 5 - France}

\scanpage{50}

we will have \deleted{that}
\[
\mu\in\widehat{\mathbb Z}^{*},\quad
\mu\equiv\deleted{\unsure{\varepsilon_3|\varepsilon_2}}(-1)^{\nu_3}(=1-2\nu_3)\quad\mathrm{mod}\ (3)
\]
and the condition on $(\mu,\varepsilon_2,\varepsilon_3)$ for being in the image of
\[
\diagram{\begin{tikzcd}
\widetilde{\underline Z}_{\rho,\sigma}\arrow[r]&\mathbf G_{\mathrm m}\times\mu_{S_0}\times\mu_{S_0}
\end{tikzcd}},
\]
which is written in general (by virtue of (62) \unsure{\&} and of (32))
\[
2\nu'-\nu_2+\nu_3\equiv0\quad(3)
\]
becomes, on multiplying by 2
\[
\underbrace{4\nu'-2\nu_2}_{\mu-1}+2\nu_3\equiv0\quad(3)
\]
i.e.
\[
\mu\equiv1-2\nu_3\quad(3),
\]
and it is automatically satisfied.

\deleted{
That is why I feel like modifying the definition of \deleted{$\unsure{Z}$} $\widetilde{\underline Z}_{\rho,\sigma}$, by defining $\underline Z_{\rho,\sigma}$ as
\[
\deleted{(63)\quad\widetilde{\underline Z}_{\rho,\sigma}=}
\left\{\begin{gathered}
\deleted{\unsure{\Gamma|\mu}}(\inserted{\mu,}\nu',\nu'',\nu_2,\nu_3,\alpha_0,\beta_0)
\in\mathbf G_{\mathrm m}\times\mathbf G_{\mathrm a}\times\unsure{\mathbf G_{\mathrm a}}\times\{0,1\}_{S_0}\times\{0,1\}_{S_0}\times\underline M\times\underline M\\
\left|\begin{gathered}
\mu=4\nu'\unsure{-|+}2\unsure{\nu_2}+1(=6\nu''-2\nu_3+1),\quad 2\nu'+\nu_3-\nu_2=3\nu'',\\
(\unsure{1}-\rho_{\underline M})(\alpha_0)-(1-\sigma_{\underline M})(\beta_0)
=\underbrace{(\deleted{\unsure{2\nu'-\nu_2+\nu_3}}\inserted{3\nu''})}_{c}\lambda_1
\end{gathered}\right.
\end{gathered}\right\}
\]
}
\margin{by \unsure{expressing|requiring} $\unsure{\mu|\nu}$\ \unsure{are written}
\[
\begin{aligned}
\mu&\equiv1-2\nu_2\quad(4)\\
\mu&\equiv1-2\nu_3\quad(6)
\end{aligned}
\unsure{\left[\vphantom{\begin{matrix}\mu\\\mu\end{matrix}}\right.}
\]
\deleted{the loop equation}
\[
\begin{aligned}
\mu&=1\unsure{-|+}2\nu_2+4\nu'&&\text{\unsure{with}}\\
\mu&=1-2\nu_3+6\nu''\\
-\nu_2+2\nu'&=-\nu_3+3\nu''&&\text{i.e.}\\
2\nu'+\nu_3-\nu_2&=3\nu''
\end{aligned}
\]}
\deleted{
\deleted{We will have here
\[
2c=6\nu'',\quad\text{and we \unsure{will|have} in fact}
\]
take $c=3\nu''$, and.}
We will have
\[
(63)\qquad
\begin{gathered}
\widetilde{\underline Z}_{\rho,\sigma}\subset\mathbf G_{\mathrm a}\times\mathbf G_{\mathrm a}\times\{0,1\}_{S_0}\times\{0,1\}_{S_0}\times\underline M\times\underline M\\
\Vert\\
\left\{(\nu',\nu'',\nu_2,\nu_3,\alpha_0,\beta_0)\,\middle|\,
\begin{aligned}
\text{a)}\quad &2\nu'\deleted{+\unsure{\nu_3}}-\nu_2=3\nu''-\nu_3\\
&\deleted{(\text{hence }4\nu'\unsure{-|+}2\nu_2+1=6\nu''-2\nu_3+1,\ \text{i.e. }\mu)}\\
&\inserted{\text{i.e. }\nu,\ \text{hence }\mu=2\nu+1}\\
\text{b)}\quad &\mu\in\Gamma\mathbf G_{\mathrm m}\\
\text{c)}\quad &\text{\unsure{loop|loop} equation}\\
&(1-\rho_{\underline M})(\alpha_0)-(1-\sigma_{\underline M})(\beta_0)=3\nu''\unsure{\underline{\lambda_1}|\lambda_1}
\end{aligned}\right\}
\end{gathered}
\]
}
\margin{49}
\margin{49\\\texttt{https://grothendieck.umontpellier.fr}\\Université de Montpellier, UMR CNRS 5149, Case courrier 051, Place Eugène Bataillon\\34095 Montpellier cedex 5 - France}

\scanpage{51}
\mspage{673}

3. The groups \deleted{$\unsure{\widetilde Z_{\rho,\sigma}}$ \unsure{revised} \illegible{2} $\mathbb O_N$} $\mathbb O_N^{*}$.

We are going to introduce an auxiliary construction,
by defining a group scheme which
corresponds to the datum of quantities
$\mu\in\Gamma\mathbf G_{\mathrm m}$, $\nu',\nu''\in\Gamma\mathbf G_{\mathrm a}$,
$\nu_2,\nu_3\in\Gamma\{0,1\}_{S_0}$, satisfying
\[
(63)\qquad
\mu=4\nu'-2\nu_2+1=6\nu''-2\nu_3+1,
\qquad 2\nu'-\nu_2=3\nu''-\nu_3.
\]

In general, let $N$\inserted{$\in\mathbb N^{*}$} be a fixed \inserted{natural} integer
(the case which interests us for what follows is
$N=12$), and let us consider the scheme \inserted{in groups, $\mathbb O_N$}
over $\mathbb Z$ associated to the extension of discrete groups
\deleted{\[(64)\qquad \mathbb O_N\]}
\[
(64)\qquad
\diagram{\begin{tikzcd}[column sep=large]
1\arrow[r]&\mathbb Z\arrow[r,"N"]&\mathbb Z\arrow[r]&\mathbb Z/N\mathbb Z\arrow[r]&1
\end{tikzcd}}
\]
which fits into an exact sequence
\[
(65)\qquad
\diagram{\begin{tikzcd}[column sep=large]
1\arrow[r]&\unsure{\mathbf E^{1}|\mathbb E^{1}}\arrow[r,"i"]&\mathbb O_N\arrow[r]&(\mathbb Z/N\mathbb Z)_{S_0}\arrow[r]&1.
\end{tikzcd}}
\]

Let us recall the \deleted{\unsure{defi}} simple construction
\[
(66)\qquad
\mathbb O_N
=(\mathbb Z)_{S_0}\mathop{\wedge}_{(N\mathbb Z)}\underline{N\mathbb Z}.
\]
\deleted{But} I say that there is \unsure{in|on} $\mathbb O_N$ a
unique Ring structure, whose
addition is the \inserted{abelian} group structure already
envisaged, and where multiplication is
obtained as follows: A section of $\mathbb O_N$ is
\deleted{local} of the form \deleted{$\unsure{k|x}$} $n+N[\lambda]$, with
the rules of calculation
\[
(67)\qquad
\left\{\begin{aligned}
\text{a)}\quad &(n+N[\lambda])+(n'+N[\lambda'])=(n+n')+N[\lambda+\lambda']\\
\text{b)}\quad &n+N[\lambda]=n'+N[\lambda']\Longleftrightarrow
\deleted{\unsure{\xi|n}}\ n+N\lambda=n'+N\lambda'\text{ in }\Gamma\mathbf G_{\mathrm a}.
\end{aligned}\right.
\]

\margin{50}
\margin{50\\\texttt{https://grothendieck.umontpellier.fr}\\Université de Montpellier, UMR CNRS 5149, Case courrier 051, Place Eugène Bataillon\\34095 Montpellier cedex 5 - France}

\scanpage{52}

We then multiply the sections by
\[
(68)\qquad
(\unsure{n}+N[\lambda])(n'+N[\lambda'])
=nn'+N[n\lambda'+n'\lambda+\lambda\lambda'],
\]
we check that it is \deleted{\unsure{biad}} compatible with
the relations (67 b), that it is biadditive,
commutative, associative, that $1+N[0]$ is
the unit\ldots{} Another way of \deleted{\unsure{it|it} \unsure{con|en}} seeing the
ring structure. \deleted{We} Let more
generally \deleted{$\mathbb O_N=\mathbb Z_{S_0}$}
$A$ be a \deleted{$\unsure{\mathbb Z|z}$} $\mathbb Z$-algebra,
\deleted{\unsure{com}} associative \inserted{\unsure{unital|unitary}} (say)
\deleted{\unsure{principal}} \inserted{equipped with an \uline{ideal}}
\margin{\inserted{two-sided}} $J$ \inserted{which is}
a free $\mathbb Z$-module of finite type, let us consider
\deleted{\unsure{an} ideal}
\deleted{the ideal $NA$} and the extension
structure
\[
(69)\qquad
\diagram{\begin{tikzcd}[column sep=large]
0\arrow[r]&\deleted{\unsure{NA}}\inserted{J}\arrow[r]&A\arrow[r]&\overbrace{A/\deleted{\unsure{NA|N}}\inserted{J}}^{B}\arrow[r]&0
\end{tikzcd}}
\]
hence a commutative group scheme
$\underline A_{[N]}$, which is obtained as
a quotient of
\deleted{$(A)_{\unsure{S|S_0}}\cdot\underline{\unsure{NB|NA}}\;\unsure{\xi\to|p}$}
semidirect \inserted{product $A_{S_0}\cdot\underline J$})
\[
(70)\qquad
\diagram{\begin{tikzcd}[column sep=large]
0\arrow[r]&\underline J\arrow[r]&A_{S_0}\times\underline J\arrow[r]&\underline A_N\arrow[r]&0\\
&\lambda\arrow[r,mapsto]&(\lambda,-\lambda)&&
\end{tikzcd}}
\]
and we find that $A_{S_0}\cdot\underline J$ has a
$\frac12$-direct product \uline{Ring} structure, and that

\margin{51}
\margin{51\\\texttt{https://grothendieck.umontpellier.fr}\\Université de Montpellier, UMR CNRS 5149, Case courrier 051, Place Eugène Bataillon\\34095 Montpellier cedex 5 - France}